\subsection{信道容量}
\subsubsection{离散信道容量}
\textbf{熵}\quad 每个发送符号的平均信息量.
\begin{equation}
    H(X)=-\sum_{i=1}^{n}P(x_i)\log_2P(x_i).
\end{equation}

\textbf{条件熵}\quad 每接收一个符号, 对于该发送符号仍剩下的平均信息量 (不确定性).
\begin{equation}
    H(X|Y)=-\sum_{j=1}^{m}P(y_j)\sum_{i=1}^{n}P(x_i|y_j)\log_2P(x_i|y_j).
\end{equation}
其中 $P(y_j)$ 表示接收符号 $y_j$ 的概率, $P(x_i|y_j)$ 表示发送 $x_i$ 接收 $y_j$ 的概率.

\textbf{互信息}\quad 平均每个符号正确传输的信息量 (每个接收符号消除发送符号的平均不确定性).
\begin{equation}
    I(X;Y)=H(X)-H(X|Y)=H(Y)-H(Y|X). \qquad \text{(b/符号)}
\end{equation}

\textbf{调整参量 (发送概率分布)}\quad 每个符号的信道容量.
\begin{equation}
    C_\text{符号}=\max_{P(x)}\{I(X;Y)\}.
\end{equation}

对于对称信道,
\begin{equation}
    C_\text{符号}=\log_2M-H(Y|X).
\end{equation}
其中, 对于对称信道, $H(Y|X)$ 恒定. 且只要 $P(X)$ 均匀分布, 互信息最大时的 $P(Y)$ 均匀分布.

若已知符号速率, 则
\begin{equation}
    C_t=rC_\text{符号}. \qquad \text{(b/s)}
\end{equation}
其中 $r$ 为每秒传输的符号数.

\subsubsection{连续信道容量}
\textbf{香农公式}\quad \textit{带限高斯白噪声}背景下连续信道容量为
\begin{equation}
    C=B\log_2\left(1+\frac{S}{n_0B}\right). \qquad \text{(b/s)}
\end{equation}
